%=====================================================================
\chapter{Paravirtualization and Device Virtualization}\label{ch:paravirt}
%=====================================================================
\locator{2}{Part II --- Hypervisors and Compute Virtualization}

\begin{outcomes}
By the end of this chapter you will be able to:
\begin{tight}
  \item \textbf{Distinguish} the three ways a guest can be given a device ---
        emulation, paravirtualization and passthrough --- by the path an I/O
        request takes in each.
  \item \textbf{Explain} the virtio split-driver model and the virtqueue, and
        \textbf{say} why batching is what makes it fast.
  \item \textbf{Describe} SR-IOV in terms of physical and virtual functions,
        and \textbf{state} what the IOMMU contributes.
  \item \textbf{Quantify} the throughput and CPU cost of each option from
        measured figures.
  \item \textbf{Choose} a device model for a stated requirement, and
        \textbf{state} the capability each choice gives up.
\end{tight}
\end{outcomes}

\begin{prereq}
\chref{architectures} (paravirtualization as a technique, and the worked
example where an emulated disk lost 81\% of its throughput) and
\chref{hardware} (VM exits, and the IOMMU without which passthrough is
unsafe). This chapter is where that 81\% is explained and removed.
\end{prereq}

\begin{keyterms}
device emulation $\bullet$ full device model $\bullet$ virtio $\bullet$ split
driver $\bullet$ front-end driver $\bullet$ back-end driver $\bullet$ virtqueue
$\bullet$ descriptor ring $\bullet$ available ring $\bullet$ used ring $\bullet$
doorbell $\bullet$ vhost $\bullet$ vhost-user $\bullet$ passthrough $\bullet$
device assignment $\bullet$ VFIO $\bullet$ SR-IOV $\bullet$ physical function
$\bullet$ virtual function $\bullet$ DPDK $\bullet$ guest additions
\end{keyterms}

\section{Three Ways to Give a Guest a Device}

A guest needs a disk and a network card. The hypervisor has one of each, or
several, and must serve many guests. There are exactly three strategies, and
they differ in how much of the real device the guest is allowed to see.

\dsfig{%
\begin{tikzpicture}[font=\scriptsize,
  lane/.style={rounded corners=3pt, line width=1pt, minimum width=4.9cm,
               minimum height=5.1cm},
  bx/.style={rounded corners=2pt, line width=0.7pt, align=center,
             minimum width=3.9cm, minimum height=0.52cm, inner sep=2pt,
             font=\scriptsize},
  gst/.style={bx, draw=secondary, fill=secondary!8, text=secondary!80!black},
  hyp/.style={bx, draw=primary, fill=primary!12, text=primary},
  dev/.style={bx, draw=neutral, fill=neutral!15, text=neutral!40!black},
  hot/.style={bx, draw=accent, fill=accent!12, text=accent!70!black},
  ar/.style={-{Stealth[length=1.7mm]}, line width=0.8pt},
  tag/.style={font=\scriptsize\bfseries, align=center}
]
% ---- emulated ----
\node[lane, draw=accent!55, fill=accent!3] at (0,0) {};
\node[tag, text=accent!70!black] at (0,2.35) {Emulated};
\node[gst] (e1) at (0,1.75) {guest driver for a\\real 1990s NIC};
\node[hot] (e2) at (0,0.95) {\textbf{VM exit} per\\register access};
\node[hyp] (e3) at (0,0.15) {hypervisor pretends\\to be that NIC};
\node[hyp] (e4) at (0,-0.65) {real driver};
\node[dev] (e5) at (0,-1.45) {hardware};
\foreach \a/\b in {e1/e2,e2/e3,e3/e4,e4/e5}{\draw[ar, accent!60] (\a)--(\b);}
\node[font=\scriptsize\itshape, text=accent!70!black, align=center] at (0,-2.15)
     {dozens of exits per I/O\\no guest changes at all};

% ---- virtio ----
\node[lane, draw=greenhl!55, fill=greenhl!3] at (5.6,0) {};
\node[tag, text=greenhl!50!black] at (5.6,2.35) {Paravirtualized (virtio)};
\node[gst] (v1) at (5.6,1.75) {guest \emph{front-end}\\driver};
\node[hot, draw=greenhl, fill=greenhl!12, text=greenhl!50!black] (v2)
     at (5.6,0.95) {shared virtqueue\\one doorbell per \emph{batch}};
\node[hyp] (v3) at (5.6,0.15) {\emph{back-end} in the\\hypervisor};
\node[hyp] (v4) at (5.6,-0.65) {real driver};
\node[dev] (v5) at (5.6,-1.45) {hardware};
\foreach \a/\b in {v1/v2,v2/v3,v3/v4,v4/v5}{\draw[ar, greenhl!60] (\a)--(\b);}
\node[font=\scriptsize\itshape, text=greenhl!50!black, align=center]
     at (5.6,-2.15) {one exit per batch of requests\\needs a driver in the guest};

% ---- passthrough ----
\node[lane, draw=purplehl!55, fill=purplehl!3] at (11.2,0) {};
\node[tag, text=purplehl!70!black] at (11.2,2.35) {Passthrough / SR-IOV};
\node[gst] (p1) at (11.2,1.75) {guest's \emph{own} driver\\for the real device};
\node[hot, draw=purplehl, fill=purplehl!12, text=purplehl!70!black] (p2)
     at (11.2,0.95) {IOMMU translates\\its DMA};
\node[dev] (p3) at (11.2,0.15) {the hardware itself};
\node[bx, draw=neutral!40, fill=white, text=neutral!50, dashed] (p4)
     at (11.2,-0.8) {hypervisor not\\on the data path};
\draw[ar, purplehl!60] (p1)--(p2); \draw[ar, purplehl!60] (p2)--(p3);
\node[font=\scriptsize\itshape, text=purplehl!70!black, align=center]
     at (11.2,-2.15) {no exits at all, near-native speed\\
                      and no live migration};
\end{tikzpicture}}%
{The same I/O request, three ways. Reading down each lane gives the cost: the
left lane crosses the boundary dozens of times per request, the middle once per
batch, the right not at all.}{device-models}

\subsection{Emulation}

The hypervisor implements, in software, a specific piece of real hardware ---
classically an Intel e1000 network card or an IDE controller --- chosen because
every operating system since the 1990s already has a driver for it. Nothing is
installed in the guest and nothing is configured. It simply works.

It is also catastrophic for throughput, and \chref{architectures} measured why.
The guest driver was written to talk to hardware, so it reads and writes device
registers, and each of those is a VM exit at several hundred cycles. One packet
costs dozens.

Emulation survives because it is the only option that needs no cooperation: an
installer booting from an unknown ISO, an appliance you cannot modify, a guest
before its drivers are installed. Every platform offers it, and no platform
expects you to keep using it.

\subsection{Paravirtualization: virtio}

If the guest can be given a driver, a far better interface is available ---
one designed for the situation rather than inherited from it.

\term{virtio} is that interface. It is a standard --- an OASIS specification,
not a vendor's product --- which is why one virtio driver in Linux works under
KVM, and under cloud hypervisors that are not KVM at all.

\begin{definitionbox}[The Split Driver Model and the Virtqueue]
A virtio device is split in two. The \term{front-end driver} runs in the guest
and presents a normal disk or network interface to the guest kernel. The
\term{back-end driver} runs in the hypervisor and does the real work. Between
them sits one or more \term{virtqueues}, in memory both can see.

\smallskip
A virtqueue is three rings in shared memory:
\begin{tight}
  \item the \term{descriptor table} --- buffers, each an address and a length;
  \item the \term{available ring} --- which descriptors the guest has offered;
  \item the \term{used ring} --- which the host has finished with.
\end{tight}

\smallskip
The guest writes buffers and advances the available ring with \emph{ordinary
memory writes}. No exit. It then rings a \term{doorbell} --- one write to one
register, one exit --- to say that work is waiting. The host takes everything
in the ring, does it, and advances the used ring.
\end{definitionbox}

\begin{conceptbox}[Batching Is the Whole Trick]
The emulated card costs an exit per register access. virtio costs an exit per
\emph{doorbell}, and one doorbell can cover a hundred queued requests.

\smallskip
Under load it gets better rather than worse, which is the opposite of the
usual pattern. A busy back-end polls the ring and finds new work already
waiting, so the guest can skip the doorbell entirely --- the specification has
an explicit flag for suppressing notifications. At high rates the exit count
per request tends toward zero.

\smallskip
That is why \chref{architectures}'s disk measurement went from 81\% lost to
4\% lost by changing one attribute. Nothing about the hardware changed; the
number of boundary crossings did.
\end{conceptbox}

\begin{conceptbox}[vhost: Moving the Back-End Out of QEMU]
In plain KVM the virtio back-end lives in QEMU, in user space. A packet
therefore crosses from the guest into the kernel, up to QEMU, back into the
kernel to reach the real driver, and out.

\smallskip
\term{vhost} moves the back-end into the host \emph{kernel}, so the data path
never enters user space; QEMU only sets the queues up. \term{vhost-user} goes
the other way for a different reason: it hands the back-end to a dedicated
user-space process --- typically a DPDK-based virtual switch --- which polls
the queues on a core of its own and never sleeps.

\smallskip
Both are the same instinct as the ladder in \chref{cpumem}: the mechanism is
settled, and all remaining performance work is removing crossings.
\end{conceptbox}

\subsection{Passthrough and SR-IOV}

The third option is to stop virtualizing the device at all and give it to one
guest. The guest loads the vendor's own driver and programs the real hardware;
the hypervisor is not on the data path.

On Linux this is done with \term{VFIO}, which unbinds the device from the host
driver and exposes it to QEMU, with the IOMMU confining its DMA to that guest's
memory. Performance is native, because nothing is in the way.

The limitation is arithmetic: one device, one guest. A host with two network
cards can serve two guests this way.

\begin{definitionbox}[SR-IOV, Physical and Virtual Functions]
\term{SR-IOV} --- single-root I/O virtualization --- is a PCIe capability that
lets one physical device present itself as many.

\smallskip
The \term{physical function} (PF) is the real device, managed by the host: it
configures the card, sets rate limits, and creates the virtual functions. Each
\term{virtual function} (VF) is a lightweight PCIe function with its own
queues, its own MAC address and its own configuration space, which can be
assigned to a guest exactly as a whole device would be.

\smallskip
The card's own hardware switch sorts traffic between the VFs. A modern server
NIC supports 64 to 128 of them, so one card serves many guests at close to line
rate, with the IOMMU keeping each VF's DMA inside its guest.
\end{definitionbox}

\begin{worked}[What Each Device Model Costs]
One 25 Gb/s network card, one guest, a small-packet forwarding test. Compute
the throughput and the cost per unit of work.

\smallskip
\textbf{Measured.}
\rowcolors{2}{white}{lightbg}
\begin{center}
\begin{tabular}{@{}L{3.5cm}rrr@{}}
\toprule
\rowcolor{primary!15}
\textbf{Model} & \textbf{Mpps} & \textbf{Host CPU, cores} &
\textbf{Guest CPU, cores} \\
\midrule
Emulated e1000  & 0.12 & 1.8 & 0.9 \\
virtio (in QEMU) & 1.40 & 1.1 & 0.8 \\
virtio + vhost  & 2.90 & 0.9 & 0.8 \\
SR-IOV VF       & 8.60 & 0.0 & 0.9 \\
Native (no guest) & 8.90 & --- & --- \\
\bottomrule
\end{tabular}
\end{center}

\smallskip
\textbf{Step 1 --- throughput against native.}
\[ \text{emulated} = \frac{0.12}{8.90} = \mathbf{1.3\%} \qquad
   \text{virtio} = \frac{1.40}{8.90} = \mathbf{15.7\%} \]
\[ \text{vhost} = \frac{2.90}{8.90} = \mathbf{32.6\%} \qquad
   \text{SR-IOV} = \frac{8.60}{8.90} = \mathbf{96.6\%} \]

\smallskip
\textbf{Step 2 --- the cost of a packet, in host cores per Mpps.}
\[ \text{emulated} = \frac{1.8}{0.12} = 15.0 \qquad
   \text{virtio} = \frac{1.1}{1.40} = 0.79 \qquad
   \text{vhost} = \frac{0.9}{2.90} = 0.31 \]
\[ \text{SR-IOV} = \frac{0.0}{8.60} = \mathbf{0} \]
Emulation is roughly \textbf{48 times} more expensive per packet than vhost,
and SR-IOV spends no host processor at all because the host is not involved.

\smallskip
\textbf{Step 3 --- the figure that decides most designs.} Look at the host-CPU
column, not the throughput column. The emulated card burns 1.8 cores to move
almost nothing; on a host of 64 cores running 30 guests, that is nearly 3\% of
the machine, per guest, spent on pretending to be a 1999 network card. The
consolidation ratio of \chref{what} is quietly destroyed by a default nobody
changed.

\smallskip
\textbf{Step 4 --- so why is SR-IOV not the answer to everything?} Because of
what it gives up, which Section~7.4 sets out: that guest can no longer be live
migrated, snapshotted with its device state, or overcommitted. Most estates
value those more than they value the last 64\% of line rate, and run virtio
with vhost. SR-IOV is for the specific guests that cannot --- a virtual
firewall, a packet broker, a trading system.
\end{worked}

\section{What Passthrough Costs You}

Passthrough is the fastest option and it is rare, because speed is not the
thing most estates are optimising.

\rowcolors{2}{white}{lightbg}
\begin{longtable}{@{}L{4.2cm}L{5.1cm}L{5.4cm}@{}}
\toprule
\rowcolor{primary!15}
\textbf{Capability} & \textbf{What happens with passthrough} &
\textbf{Why} \\
\midrule
\endhead
Live migration (\chref{migration}) & \textbf{Lost}, or needs an elaborate
bonded-failover arrangement & The device's internal state is in silicon the
hypervisor cannot read, so there is nothing to copy to the other host \\
Memory overcommit (\chref{cpumem}) & \textbf{Lost} for that guest & The device
performs DMA to guest memory at any moment, so every page must be pinned in
real memory and cannot be ballooned or swapped \\
Snapshots & Lost or incomplete & The device's state is not part of the guest's
memory image \\
Density & One VF per guest, 64--128 per card & A hard hardware limit, where
virtio has none \\
Hardware independence & Lost & The guest now has a driver for a specific card.
Move it to a host with a different NIC and it will not work \\
Host visibility & Reduced & The host cannot see or shape that traffic; firewall
and monitoring must move into the guest or onto the card \\
\bottomrule
\end{longtable}

\begin{alertbox}[Memory pinning is the consequence people forget]
Live migration is the loss everyone cites. Pinning is the one that bites
quietly.

\smallskip
A device doing DMA writes to host physical addresses it was given earlier. If
the hypervisor moved, compressed, ballooned or swapped that page in the
meantime, the device writes over whatever is there now. So every page of a
passthrough guest's memory must be locked in RAM for the guest's whole life.

\smallskip
A 64 GB guest with an assigned device consumes 64 GB of the host whether it
uses it or not --- and the entire reclamation ladder of \chref{cpumem} is
unavailable for it. On a host sized with overcommit in mind, two or three such
guests change the arithmetic completely.
\end{alertbox}

\begin{pitfall}
\begin{tight}
  \item \textbf{Leaving guests on emulated devices and concluding
        virtualization is slow.} It is the commonest cause of a bad first
        impression, and it is one attribute in the guest's configuration.
  \item \textbf{Installing guest tools ``for the clipboard''.} The drivers are
        the point. The clipboard is a side effect.
  \item \textbf{Assigning a device without checking IOMMU groups.} Devices
        share groups, and the whole group must be assigned together. Passing
        one function of a multi-function card may drag its siblings --- which
        might include a device the host needs.
  \item \textbf{Expecting to migrate a guest with an assigned device.} Plan the
        exception before you build it, not when you first need to patch the
        host.
  \item \textbf{Forgetting that SR-IOV traffic bypasses the host switch.}
        Packets between two VFs on the same card never reach the host, so host
        firewall rules, flow logging and traffic shaping do not see them.
        \chref{network} and \chref{security} both depend on this.
  \item \textbf{Treating virtio as a Linux feature.} It is an open
        specification with signed Windows drivers (virtio-win) that are free
        and are what every cloud Windows image already uses.
\end{tight}
\end{pitfall}

\begin{lab}[The Same Card, Three Ways]
Measures the throughput and the host CPU cost of an emulated NIC, a virtio NIC
and virtio with vhost. SR-IOV is included as a read-only step because it needs
a server NIC. Expect forty minutes.

\begin{lstlisting}[language=bash]
#!/usr/bin/env bash
# ch07_devices.sh -- emulation, virtio and vhost, measured.
set -euo pipefail
ip_of () {
  sudo virsh domifaddr "$1" | awk '/ipv4/{split($4,a,"/"); print a[1]}'
}
G () { ssh -o StrictHostKeyChecking=no "ubuntu@$(ip_of vss1)" "$1"; }

# --- 1. What device model is the guest using right now? ----------------
sudo virsh dumpxml vss1 | sed -n '/<interface/,/<\/interface>/p'
# <model type='virtio'/> or type='e1000'. This single line is the chapter.

# --- 2. Set up the measurement: iperf3 on the host, client in guest. ---
sudo apt-get install -y iperf3 sysstat
iperf3 -s -D                       # server on the host, port 5201
sudo virsh start vss1 || true ; sleep 25
G "sudo apt-get -qq install -y iperf3"

# --- 3. Emulated e1000: a software imitation of a 1999 card. ----------
sudo virsh destroy vss1
sudo virt-xml vss1 --edit --network model=e1000
sudo virsh start vss1 ; sleep 25
mpstat 1 12 > host-e1000.txt &     # host CPU while the test runs
G "iperf3 -c 192.168.122.1 -t 10 -P 4" | tail -4
wait ; tail -2 host-e1000.txt

# --- 4. virtio: the split-driver model of section 7.2. ----------------
sudo virsh destroy vss1
sudo virt-xml vss1 --edit --network model=virtio
sudo virsh start vss1 ; sleep 25
mpstat 1 12 > host-virtio.txt &
G "iperf3 -c 192.168.122.1 -t 10 -P 4" | tail -4
wait ; tail -2 host-virtio.txt

# --- 5. virtio + vhost: back-end moved into the host kernel. ----------
# vhost_net is usually already on; confirm, and confirm the guest uses it.
lsmod | grep vhost_net || sudo modprobe vhost_net
sudo virsh destroy vss1
sudo virt-xml vss1 --edit --network driver.name=vhost
sudo virsh start vss1 ; sleep 25
mpstat 1 12 > host-vhost.txt &
G "iperf3 -c 192.168.122.1 -t 10 -P 4" | tail -4
wait ; tail -2 host-vhost.txt

# --- 6. Look at the virtqueues themselves. ----------------------------
# One queue pair per vCPU is the usual arrangement; each is a ring in
# memory that both sides can see, exactly as the definition box says.
G "ls /sys/class/net/*/queues/ ; ethtool -l ens3 2>/dev/null || true"

# --- 7. SR-IOV, if the host has a server NIC that supports it. --------
# Read-only: creating VFs needs a capable card and IOMMU enabled.
for d in /sys/class/net/*/device/sriov_totalvfs; do
  [ -e "$d" ] && echo "$d -> $(cat "$d") VFs available"
done
ls /sys/kernel/iommu_groups 2>/dev/null | wc -l

# --- 8. Restore and tear down. ----------------------------------------
sudo virsh destroy vss1 || true
pkill iperf3 || true
\end{lstlisting}

\textbf{What to notice.} The three throughput figures should differ by roughly
the ratios in the worked example, but look at the \texttt{mpstat} output rather
than the bandwidth: the emulated run burns host processor to move very little,
and that --- not the bandwidth --- is what decides how many guests the host can
hold.
\end{lab}

\begin{checkpoint}
\begin{tightnum}
  \item A guest moves 0.2 Mpps while the host spends 1.6 cores serving it.
        Give the cost in host cores per Mpps, say which device model this
        almost certainly is, and give the one change you would make.
  \item Explain why a virtio guest's exit rate per request \emph{falls} as load
        rises, and name the mechanism.
  \item A team wants SR-IOV for twelve guests on a host that must stay
        patchable without downtime. Name the conflict, and give two ways to
        resolve it.
\end{tightnum}
\end{checkpoint}

\begin{chaptersummary}
\begin{tight}
  \item Three device models: emulation (no guest changes, dozens of exits per
        I/O), paravirtualization (a driver in the guest, one exit per batch),
        and passthrough (the guest drives real hardware, no exits).
  \item virtio splits the driver in two around a virtqueue in shared memory.
        The guest fills the ring with ordinary writes and rings a doorbell once
        per batch, so the crossing rate per request falls as load rises.
  \item vhost moves the back-end into the host kernel; vhost-user moves it to a
        polling user-space switch. Both remove crossings rather than making
        them cheaper.
  \item SR-IOV makes one card present many virtual functions, each assignable
        to a guest, with the IOMMU confining each one's DMA.
  \item Passthrough costs live migration, memory overcommit, snapshots,
        hardware independence and host visibility. Memory must be pinned for
        the guest's whole life.
  \item Most estates run virtio with vhost, and reserve passthrough for the few
        guests that genuinely need line rate.
\end{tight}
\end{chaptersummary}

\begin{reviewq}
  \item \textbf{(MCQ)} A virtqueue consists of: (a)~a single FIFO buffer;
        (b)~a descriptor table, an available ring and a used ring;
        (c)~one ring per device register; (d)~a copy of the guest's page
        tables.
  \item \textbf{(MCQ)} Assigning a physical device to a guest requires every
        page of that guest's memory to be pinned because: (a)~the guest runs
        faster; (b)~the device performs DMA to addresses that must stay valid;
        (c)~the IOMMU has no page tables; (d)~live migration needs it.
  \item \textbf{(Short)} Explain the difference between a physical function and
        a virtual function, and say what each is used for.
  \item \textbf{(Short)} Why does emulation still exist, given that it is
        roughly fifty times more expensive per packet than the alternatives?
  \item \textbf{(Short)} State three capabilities a guest loses when a device
        is assigned to it, with the reason for each.
  \item \textbf{(Applied)} A virtual firewall must inspect 10 Gb/s of
        small-packet traffic. The estate's standard is virtio with vhost and
        all guests are live migrated for monthly patching. Recommend a device
        model, state what it costs, and describe how you would keep the host
        patchable anyway.
  \item \textbf{(Applied)} A host of 64 cores runs 30 guests, all on emulated
        e1000 adapters because that was the template default. Using the worked
        example's figures, estimate the host processor currently spent on
        device emulation, say what you would change, and describe how you would
        roll the change out without an outage.
\end{reviewq}
